\chapter{Cơ sở lý thuyết và ánh xạ kiến thức}

\section{Biểu diễn ảnh và không gian màu}
Phần này cần trình bày các khái niệm nền tảng như ảnh số, pixel, độ phân giải, kênh màu và các không gian màu thường dùng như RGB, HSV hoặc Lab. Đây là cơ sở để mô tả cách dữ liệu ảnh được lưu trữ và xử lý trước khi đưa vào mô hình.

\section{Lọc ảnh và khử nhiễu}
Nội dung ở đây tập trung vào các bộ lọc thường dùng như Gaussian, median hoặc bilateral filtering. Mục tiêu là giảm nhiễu, làm mượt ảnh hoặc bảo toàn biên tùy theo yêu cầu bài toán.

\section{Phát hiện biên, ngưỡng hóa và hình thái học}
Phần này mô tả các kỹ thuật như Sobel, Canny, thresholding, erosion, dilation, opening, closing và hole filling. Những phương pháp này thường dùng để tách tiền cảnh, làm sạch mask và hỗ trợ phân đoạn.

\section{Phân đoạn, thành phần liên thông và đặc trưng}
Đây là nơi trình bày các kỹ thuật như connected components, contour analysis, area measurement, shape analysis, keypoints, descriptors và matching. Nếu đề tài có phân đoạn hoặc nhận dạng theo mẫu, phần này là trọng tâm.

\section{Phân loại, nhận dạng và đánh giá}
Chương này cần giải thích các mô hình dùng cho phân loại hoặc nhận dạng, chẳng hạn SVM, KNN, Random Forest hoặc một mô hình CNN nhẹ. Đồng thời cần nêu rõ các chỉ số đánh giá phù hợp như Accuracy, Precision, Recall, F1, IoU, Dice, mAP hoặc runtime.

\section{Ánh xạ giữa kiến thức môn học và nhiệm vụ đề tài}
Bảng~\ref{tab:course_mapping} tóm tắt cách các mảng kiến thức trong học phần được đưa vào bài toán thực tế của báo cáo.

\begin{table}[H]
    \centering
    \caption{Ánh xạ giữa kiến thức môn học và các nhiệm vụ thường gặp trong đề tài thị giác máy tính.}
    \label{tab:course_mapping}
    \resizebox{\textwidth}{!}{%
    \begin{tabular}{|p{0.24\textwidth}|p{0.38\textwidth}|p{0.30\textwidth}|}
        \hline
        \textbf{Kiến thức} & \textbf{Nhiệm vụ trong đề tài} & \textbf{Ví dụ/ghi chú} \\
        \hline
        Biểu diễn ảnh, kênh màu, không gian màu & Đọc ảnh, chuyển đổi RGB/HSV/Lab, chuẩn hóa màu, phân tích histogram & Tùy bài toán nhận dạng, phân đoạn hoặc theo dõi đối tượng \\
        \hline
        Lọc ảnh và khử nhiễu & Gaussian/median/bilateral filtering, làm mịn trước khi trích xuất đặc trưng & Hữu ích khi ảnh có nhiễu hoặc nền phức tạp \\
        \hline
        Phát hiện biên và gradient & Sobel, Canny, phát hiện cạnh, bản đồ gradient & Phục vụ trích vùng, đo hình dạng, hoặc tiền xử lý cho segmentation \\
        \hline
        Ngưỡng hóa và hình thái học & Thresholding, erosion, dilation, opening, closing, hole filling & Tách tiền cảnh, làm sạch mask, nối vùng đứt gãy \\
        \hline
        Thành phần liên thông và contour & Đếm đối tượng, đo diện tích, chu vi, bounding box, shape analysis & Hữu ích cho nhận dạng vật thể đơn giản hoặc đo đạc \\
        \hline
        Trích chọn và khớp đặc trưng & Keypoints, descriptors, matching, homography & Dùng trong stitching, tracking hoặc nhận dạng theo mẫu \\
        \hline
        Phân cụm và phân đoạn & K-means, watershed, superpixel, semantic/instance segmentation & Phân vùng đối tượng hoặc khu vực quan tâm \\
        \hline
        Thiết kế đặc trưng & Histogram, texture, shape, LBP/GLCM/HOG & Đầu vào cho mô hình học máy truyền thống \\
        \hline
        Phân loại và nhận dạng & SVM, KNN, Random Forest, lightweight CNN & So sánh baseline và phương pháp đề xuất \\
        \hline
        Đánh giá hệ thống & Accuracy, precision-recall, IoU, Dice, runtime & Báo cáo cả định lượng lẫn định tính \\
        \hline
    \end{tabular}%
    }
\end{table}

\section{Các chỉ số đánh giá}
\begin{itemize}
    \item Với bài toán phân loại, cần nêu Accuracy, Precision, Recall và F1-score.
    \item Với bài toán phân đoạn, cần nêu IoU, Dice và mIoU.
    \item Với bài toán phát hiện đối tượng, cần nêu Precision-Recall, mAP và số lỗi phát hiện.
    \item Nếu đề tài có yêu cầu triển khai, nên bổ sung runtime, bộ nhớ và độ ổn định của hệ thống.
\end{itemize}

